\setcounterpageref{ex}{0}
\TaoPhieuKiemTra{10}

\setcounter{ex}{0}
\setcounter{bt}{0}

\textbf{\large I. PHẦN TRẮC NGHIỆM (3,0 điểm)}

{\itshape Khoanh tròn vào chữ cái đặt trước câu trả lời đúng. Mỗi câu đúng được 0,5 điểm.}

\begin{ex} % Câu 1
Số liền trước của số lớn nhất có hai chữ số là
\choice
{$99$}
{$100$}
{\True $98$}
{$97$}
\loigiai{
Số lớn nhất có hai chữ số là $99$, số liền trước là $98$.\\
\textbf{Đáp án đúng: C}
}
\end{ex}

\begin{ex} % Câu 2
Trong phép tính $23 + 51 = 74$, kết quả phép tính được gọi là
\choice
{Hiệu}
{\True Tổng}
{Số hạng}
{Số bị trừ}
\loigiai{
Kết quả phép cộng gọi là Tổng.\\
\textbf{Đáp án đúng: B}
}
\end{ex}

\begin{ex} % Câu 3
Con mèo nhà em cân nặng khoảng
\choice
{$3\text{ cm}$}
{$3\text{ }l$}
{$3\text{ dm}$}
{\True $3\text{ kg}$}
\loigiai{
Khối lượng con mèo đo bằng ki-lô-gam ($3\text{ kg}$).\\
\textbf{Đáp án đúng: D}
}
\end{ex}

\begin{ex} % Câu 4
Hình vuông có bao nhiêu cạnh bằng nhau?
\choice
{\True $4$ cạnh}
{$2$ cạnh}
{$3$ cạnh}
{$5$ cạnh}
\loigiai{
Hình vuông có $4$ cạnh bằng nhau.\\
\textbf{Đáp án đúng: A}
}
\end{ex}

\begin{ex} % Câu 5
Điền dấu thích hợp: $9 + 4 \dots 8 + 5$.
\choice
{$>$}
{\True $=$}
{$<$}
{Không so sánh được}
\loigiai{
$9 + 4 = 13$; $8 + 5 = 13$. Vì $13 = 13$ nên điền dấu $=$.\\
\textbf{Đáp án đúng: B}
}
\end{ex}

\begin{ex} % Câu 6
Tính: $16 - 7 + 5 = \dots$
\choice
{$13$}
{$15$}
{\True $14$}
{$12$}
\loigiai{
$16 - 7 = 9$; $9 + 5 = 14$.\\
\textbf{Đáp án đúng: C}
}
\end{ex}

\vspace{4mm}
\textbf{\large II. PHẦN TỰ LUẬN (7,0 điểm)}

{\itshape Thí sinh trình bày chi tiết bài làm vào giấy kiểm tra.}

\begin{bt}[1,0 điểm]
Tính nhẩm:
\begin{multicols}{2}
\begin{enumerate}[a)]
  \item $8 + 4$
  \item $7 + 5$
  \item $12 - 7$
  \item $14 - 8$
\end{enumerate}
\end{multicols}
\loigiai{
\begin{itemize}
  \item $8 + 4 = 12$
  \item $7 + 5 = 12$
  \item $12 - 7 = 5$
  \item $14 - 8 = 6$
\end{itemize}
}
\end{bt}

\begin{bt}[1,5 điểm]
Đặt tính rồi tính:
\begin{multicols}{3}
\begin{enumerate}[a)]
  \item $43 + 37$
  \item $75 - 31$
  \item $61 + 56$
\end{enumerate}
\end{multicols}
\loigiai{
\begin{itemize}
  \item $43 + 37 = 80$
  \item $75 - 31 = 44$
  \item $61 + 56 = 117$
\end{itemize}
}
\end{bt}

\begin{bt}[1,0 điểm]
Tính giá trị biểu thức kèm đơn vị đo:
\begin{enumerate}[a)]
  \item $26\text{ kg} + 9\text{ kg} - 10\text{ kg}$
  \item $35\text{ }l - 11\text{ }l + 8\text{ }l$
\end{enumerate}
\loigiai{
\begin{enumerate}[a)]
  \item $26\text{ kg} + 9\text{ kg} - 10\text{ kg} = 35\text{ kg} - 10\text{ kg} = 25\text{ kg}$.
  \item $35\text{ }l - 11\text{ }l + 8\text{ }l = 24\text{ }l + 8\text{ }l = 32\text{ }l$.
\end{enumerate}
}
\end{bt}

\begin{bt}[1,5 điểm]
Bạn Vy có $25$ bức tranh. Bạn Bảo có ít hơn bạn Vy $8$ bức tranh. Hỏi bạn Bảo có bao nhiêu bức tranh?
\loigiai{
Số bức tranh bạn Bảo có là:
\[ 25 - 8 = 17 \text{ (bức tranh)} \]
\textbf{Đáp số:} $17$ bức tranh.
}
\end{bt}

\begin{bt}[1,0 điểm]
Một thùng chứa đầy $45\text{ }\text{kg}$ gạo nếp. Buổi sáng cửa hàng đã bán ra $14\text{ }\text{kg}$. Hỏi trong thùng còn lại bao nhiêu ki-lô-gam gạo nếp?
\loigiai{
Trong thùng còn lại số ki-lô-gam gạo nếp là:
\[ 45 - 14 = 31\text{ }\text{kg} \]
\textbf{Đáp số:} $31\text{ }\text{kg}$ gạo nếp.
}
\end{bt}

\begin{bt}[1,0 điểm]
Điền tiếp hai số thích hợp vào chỗ chấm để hoàn thành quy luật dãy số: $3, 6, 9, 12, \dots, \dots$
\loigiai{
Quy luật: Mỗi số sau bằng số trước cộng thêm $3$.\\
Số tiếp theo: $12 + 3 = 15$.\\
Số tiếp theo nữa: $15 + 3 = 18$.\\
\textbf{Đáp số:} $15$ và $18$.
}
\end{bt}
